\section{Limitations}

The study uses a purposive single-coder corpus from one public specialist forum. Selection was designed for conceptual coverage and difficult cases, not representative sampling. Discussion recurrence cannot estimate field-wide incidence, product reliability, or the distribution of maintenance work. Public participation also reflects product installed base, system complexity, user willingness to post, and community history.

The public record is incomplete. Support can migrate into private channels, and technical or scientific constraints can prevent disclosure. We treat migration as an analytical outcome and cannot reconstruct the hidden work. The corpus therefore provides a view of public articulation and repair, not the complete cooperative system surrounding every case.

The coding model is one interpretation of the material. Explicit boundary rules, hard-case adjudication, counterexamples, source audits, and executable coherence checks strengthen transparency without replacing independent qualitative replication. A blind hard-case set is available for a future second pass. The present paper reports no agreement statistic.

The design implications are derived from observed mechanisms and released as proposals. The troubleshooting form, knowledge record, interface registry, and event schema have not yet been evaluated as interventions. Their effects on participation, resolution, privacy, repair, and maintenance require prospective studies with laboratories and community members.

Finally, the analysis focuses on laboratory automation, where software and physical state are tightly coupled. Some findings should transfer to other cyber-physical technical communities, although local units, risks, institutional arrangements, and evidence practices require separate study.
